\documentclass[10pt,a4paper]{amsart}
\usepackage[margin=19mm]{geometry}
\usepackage{amsmath,amssymb,mathtools}
\usepackage[scr=rsfs,cal=euler]{mathalfa}
\usepackage{xfrac}
\usepackage[unicode,hidelinks,bookmarks=false]{hyperref}
\newcommand{\ie}{i.e.\ }
\newcommand{\cf}{cf.\ }
\newcommand{\resp}{respectively\ }
\newcommand{\Subsection}{\subsection}
\providecommand{\nobreakdash}{\nobreak\hbox{-}}
\setlength{\emergencystretch}{2em}
\multlinegap0pt
\usepackage{amssymb}
\usepackage{mathtools}
\newtheorem{theorem}{Theorem}
\title{Spectral Codes: A Geometric Formalism for Quantum Error Correction}
\author{Satoshi Kanno, Yoshi-aki Shimada}
\date{Source: arXiv:2601.19765v1 (CC BY 4.0)}
\begin{document}
\maketitle

\noindent\emph{Source and licence.} Spectral Codes: A Geometric Formalism for Quantum Error Correction. Version arXiv:2601.19765v1 (CC BY 4.0), distributed by arXiv under the Creative Commons Attribution 4.0 International licence, http://creativecommons.org/licenses/by/4.0/, as recorded in the arXiv licence metadata for this exact version and shown on the abstract page https://arxiv.org/abs/2601.19765v1. Full licence notice: \texttt{licenses/arxiv-2601.19765-CC-BY-4.0.txt}.\par\smallskip

\noindent\textit{Editorial scope.} Counted unit: the article's theorem theorem (source0.tex lines 2010--2018). The article's complete original proof of it is delivered with it (source0.tex lines 2020--2040, one \texttt{proof} environment of 21 lines). No mathematics is written, repaired or re-derived: the statement and the proof are the article's own lines, in the article's own notation, and no retained line is substituted or re-flowed.  No further source-local context is needed: every cross-reference of every retained line resolves inside the two delivered spans.  Nothing was omitted from any retained span, so the package carries no omission record.  No retained line contains a citation, so the wrapper carries no bibliography and no bibliography entry.  No retained line contains a figure environment, an image-inclusion command or drawing code, so no image is delivered and the package declares no asset dependency.  Editorial formatting only, in addition: an AMS \texttt{amsart} preamble that restates the article's own declarations and macros used by the retained lines, namely the environments \texttt{theorem} on separate counters, together with the standard packages those lines need.  The retained environments print in this document as Theorem 1 in order of appearance, whereas the article prints the same items with the numbering of its own full document.  The complete native source of the article as distributed in the arXiv e-print for this exact version is bundled unchanged as \texttt{sources/193484/source0.tex}.
\begin{theorem}[Distance of the topological code]
The code distance is given by
\[
    d_D(P)
    =
    \min\{\mathrm{wt}(F_\rho)\mid
    \rho\ \text{is a non-contractible closed ribbon}\}.
\]
\end{theorem}

\begin{proof}
We consider closed ribbon operators.

\emph{(Contractible case).}
If $\rho$ is contractible, then the corresponding operator $F_\rho$ can be written as a product of vertex and plaquette operators.
Hence, on the code space,
\[
    P F_\rho P = \lambda_\rho P
\]
for some scalar $\lambda_\rho$, and $F_\rho$ acts trivially on $\mathcal C$.

\emph{(Non-contractible case).}
If $\rho$ is non-contractible, then $F_\rho$ cannot be generated by local stabilizers.
It acts as a nontrivial logical operator on $\mathcal C$, and therefore
\[
    P F_\rho P \notin \mathbb{C}P.
\]

By definition, the distance of the code is the minimal weight of an operator acting nontrivially on the code space.
Since only non-contractible closed ribbons act nontrivially, and their weights are measured by $\mathrm{wt}(F_\rho)$, the stated formula follows.
\end{proof}

\end{document}
